\documentclass[12pt,a4paper]{article}

\usepackage[T1]{fontenc}
\usepackage[utf8]{inputenc}
\usepackage[polish]{babel}
\usepackage{lmodern}
\usepackage[a4paper,margin=2.2cm]{geometry}
\usepackage{graphicx}
\usepackage{enumitem}
\usepackage{tabularx}
\usepackage{booktabs}
\usepackage{fancyhdr}
\usepackage{hyperref}

\hypersetup{
    colorlinks=false,
    pdftitle={Noc Politechniki Wrocławskiej 2026},
    pdfauthor={Jarosław Rybak}
}

\pagestyle{fancy}
\fancyhf{}
\lhead{Noc Politechniki Wrocławskiej 2026}
\rhead{Jarosław Rybak}
\cfoot{\thepage}

\setlist[itemize]{leftmargin=1.2cm,itemsep=2pt,topsep=3pt}
\setlist[enumerate]{leftmargin=1.2cm,itemsep=2pt,topsep=3pt}

\newcommand{\sectiontitle}[1]{%
    \section*{#1}
}

\begin{document}

\begin{titlepage}
    \centering
    \vspace*{2.1cm}
    {\Huge\bfseries Noc Politechniki Wrocławskiej 2026\par}
    \vspace{1.8cm}
    {\Large Projekt interaktywnej konsoli do gier\par}
    \vspace{2.2cm}
    {\Large\bfseries Autor: Jarosław Rybak\par}
    \vspace{0.3cm}
    {\large Politechnika Wrocławska\par}
    \vspace{1.8cm}
    \begin{minipage}{0.78\textwidth}
        \centering
        \small
        Projekt to mała konsola do gier z ekranem, joystickiem, diodami LED
        i czujnikiem ruchu. Powstała jako ciekawy pokaz elektroniki,
        programowania i działania urządzeń podczas Nocy Politechniki
        Wrocławskiej 2026.

        \vspace{0.5cm}

        Użytkownik może poruszać się po menu, wybierać gry, sterować
        klasycznymi trybami joystickiem oraz przechylać płytkę, aby
        sterować kropką za pomocą czujnika MPU6050. Na osobnym ekranie
        można obserwować dane czujnika w czasie rzeczywistym.
    \end{minipage}
    \vspace{1.3cm}

    \begin{minipage}{0.78\textwidth}
        \textbf{Projekt zawiera:}
        \begin{itemize}
            \item kolorowe menu i kilka gier zręcznościowych,
            \item ekran TFT o rozdzielczości 160$\times$128 pikseli,
            \item joystick analogowy i przycisk,
            \item pięć diod LED pokazujących liczbę żyć,
            \item czujnik ruchu MPU6050 z akcelerometrem i żyroskopem,
            \item ekran telemetrii pokazujący działanie czujnika,
            \item tryby PAINT i BALANS sterowane przechyleniem płytki.
        \end{itemize}
    \end{minipage}
    \vfill
    {\large 2026\par}
\end{titlepage}

\newpage
\sectiontitle{Opcje menu i gry}

\begin{itemize}
    \item \textbf{SNAKE} -- klasyczny wąż porusza się po planszy z widoczną
    kratką. Zjedzenie jedzenia zwiększa wynik, a okresowo pojawia się bonus
    punktowy. Kolizja ze ścianą lub własnym ciałem kończy próbę.
    \item \textbf{KROPKA} -- żółta kropka unika czerwonych bomb. Przeszkody
    nadlatują z góry, dołu, lewej i prawej strony, a ich liczba rośnie
    wraz z czasem.
    \item \textbf{LABIRYNT} -- gracz prowadzi kropkę od punktu A do punktu B,
    omijając czerwone ściany. Dostępnych jest dziesięć poziomów trudności,
    z coraz mniejszym celem i krótszym limitem czasu.
    \item \textbf{CZUJNIKI} -- osobne podmenu zawierające telemetrię,
    sensorowy tryb PAINT, grę BALANS i test pięciu diod LED.
\end{itemize}

\subsection*{Tryby wykorzystujące czujnik ruchu}

\textbf{PAINT} pozwala malować po ekranie przez przechylanie płytki.
Na początku program zapamiętuje pozycję płaską. Dzięki temu kropka nie
drży, kiedy urządzenie leży na stole. Obrót płytki może zmieniać kolor.

\textbf{BALANS} to prosta próba zręczności. Zielony cel zmienia położenie,
a punkty dostaje się za utrzymanie żółtej kropki w jego obrębie. Gra korzysta
tylko z danych czujnika ruchu.

\newpage

\sectiontitle{Komponenty i połączenia}

\begin{tabularx}{\textwidth}{@{}lX@{}}
\toprule
\textbf{Element} & \textbf{Zastosowanie} \\
\midrule
NUCLEO-L476RG / STM32L476RG & Mikrokontroler wykonujący program konsoli \\
ST7735S TFT, 160$\times$128 & Menu, gry, telemetria i wizualizacja przechyłu \\
Analogowy joystick & Wybór pozycji, sterowanie klasycznymi grami i przycisk akcji \\
Crowtail-MPU6050 & Akcelerometr i żyroskop do telemetrii oraz sterowania ruchem \\
KAmodVL6180 & Czujnik odległości ToF, przygotowany do przyszłych gestów \\
5 diod LED + rezystory & Wizualizacja liczby żyć i ręczny test wyjść GPIO \\
Płytka stykowa i przewody & Wspólne zasilanie oraz porządkowanie połączeń \\
\bottomrule
\end{tabularx}

\subsection*{Najważniejsze piny}

\begin{tabularx}{\textwidth}{@{}l l X@{}}
\toprule
\textbf{Element} & \textbf{Pin STM32} & \textbf{Funkcja} \\
\midrule
Joystick VRx & PA0 & ADC1\_IN5, oś pozioma \\
Joystick VRy & PA1 & ADC1\_IN6, oś pionowa \\
Joystick SW & PB0 & Przycisk z wejściem pull-up \\
MPU6050 i VL6180 & PB8 / PB9 & Wspólna magistrala I2C: SCL / SDA \\
TFT MOSI / SCK & PC3 / PB10 & SPI2, dane i zegar \\
TFT DC / CS / RST & PB11 / PB12 / PB2 & Sterowanie wyświetlaczem \\
LED-y życia & PC10, PC9, PC8, PC7, PC6 & Wyjścia GPIO, osobny rezystor dla każdej LED \\
\bottomrule
\end{tabularx}

Wszystkie moduły pracują ze wspólną masą. Układ jest zasilany z napięcia
3,3 V, a każda dioda LED ma osobny rezystor szeregowy. MPU6050 i VL6180
mogą współdzielić linie I2C, ponieważ mają różne typowe adresy:
\texttt{0x68} oraz \texttt{0x29}.

\sectiontitle{Jak elementy komunikują się ze sobą?}

Najprościej można powiedzieć, że każdy pin jest małym połączeniem, przez
które płynie informacja albo zasilanie. Pin może podawać stan niski lub
wysoki, czyli w praktyce sygnał podobny do zera i jedynki. Dzięki temu
mikrokontroler może włączyć diodę, sprawdzić przycisk albo przesłać dane
do ekranu. Niektóre piny mają także dodatkowe zadania, na przykład mogą
odczytywać napięcie z joysticka.

\subsection*{Obraz na ekranie}

Ekran TFT nie może otrzymać przypadkowego ciągu szybkich sygnałów. Musi
wiedzieć, które dane są właśnie przesyłane i kiedy ma je odczytać. W tym
projekcie obraz jest wysyłany przez SPI. Jedna linia przekazuje dane, druga
ustala rytm przesyłania, a kolejne sygnały informują, do którego urządzenia
trafiają dane i czy są one poleceniem, czy fragmentem obrazu. Ten rytm
nazywa się zegarem. Można to porównać do podawania klocków w równych
odstępach: ekran odbiera każdy klocek w odpowiednim momencie i składa z nich
kolorowe piksele.

\subsection*{Joystick i diody}

Joystick ma dwie osie. Ich położenie zmienia napięcie, a mikrokontroler
mierzy je za pomocą ADC. ADC zamienia napięcie na liczbę, dzięki czemu
program rozpoznaje ruch w lewo, prawo, górę lub dół. Przycisk joysticka
działa prościej: po naciśnięciu zmienia stan jednego pinu. Diody LED są
sterowane sygnałami wyjściowymi. Program podaje stan wysoki, aby je zapalić,
albo niski, aby je zgasić. Każda dioda ma własny rezystor, który ogranicza
prąd i chroni zarówno diodę, jak i pin mikrokontrolera.

\subsection*{Czujniki na wspólnej magistrali}

MPU6050 i VL6180 korzystają z I2C. W tej metodzie dwa przewody są wspólne
dla wielu urządzeń: jeden przenosi dane, a drugi pomaga ustalić ich rytm.
Każdy czujnik ma własny adres, dlatego mikrokontroler może powiedzieć:
„teraz odpowiada MPU6050”, a później „teraz odpowiada VL6180”. To pozwala
podłączyć oba czujniki równolegle bez używania osobnych przewodów dla
każdego z nich.

Ważne jest również połączenie wspólnej masy. Wszystkie elementy muszą mieć
wspólny punkt odniesienia, inaczej sygnały mogą być dla siebie
niezrozumiałe. Dlatego w projekcie przewody GND modułów są połączone razem,
a elementy są zasilane napięciem zgodnym z ich dokumentacją.

\sectiontitle{Platforma i ciekawostki techniczne}

Projekt został wykonany w programie \textbf{STM32CubeIDE 2.2.0}. Kod jest
napisany w języku C dla płytki NUCLEO-L476RG z układem STM32L476RG.
STM32 to mały komputer na jednej płytce, który wykonuje cały program.
Ekran komunikuje się z nim przez SPI, joystick przez ADC, a czujniki przez
I2C. Są to po prostu różne sposoby przesyłania informacji między elementami.

\begin{itemize}
    \item MPU6050 jest automatycznie wykrywany pod typowym adresem $0x68$
    lub alternatywnym $0x69$.
    \item Kalibracja zapisuje średnią z 5 próbek przy płasko położonej
    płytce i traktuje ją jako pozycję neutralną.
    \item Odczyty czujnika są wygładzane, dzięki czemu kropka nie drży
    przy małych, przypadkowych ruchach.
    \item Telemetria odświeża tylko zmieniające się pola liczbowe i marker
    przechyłu, zamiast rysować cały ekran od początku.
    \item Pięć LED-ów pokazuje życia w kolejności: czerwona, czerwona,
    żółta, zielona, zielona.
    \item Generator pseudolosowy tworzy pozycje jedzenia, bomb i celów bez
    dodatkowych bibliotek.
\end{itemize}

\subsection*{Po co powstał projekt?}

Projekt pokazuje, jak połączyć elektronikę z programowaniem. Osoba oglądająca
może przechylić płytkę, zobaczyć zmianę danych na ekranie i od razu użyć
tego ruchu w grze. Dzięki temu łatwiej zrozumieć, czym są czujniki,
mikrokontrolery i programowanie urządzeń.

\end{document}
